\originalsection{18.212 Spring 2019：课程作业}
本节选用 Alexander Postnikov 的三份课程作业中的19道原题. 官方课堂解答由 Andrew Lin 记录, 其中只有部分题有解答, 而且若干构造仅给出提示. 每题注明原题号及答案参考, 本稿补足条件、算法及证明; 没有官方答案的题明确标为编者独立补解. 这里的题号不沿用官方解答文件重新编排的局部编号.

\par\medskip\noindent\textbf{Problem Set 1：随机游走、偏序与包含双射}

\begin{exercise}\label{mit:18212-pset1-q2}
\textbf{18.212 Spring 2019, Problem Set 1, Problem 2.}
在整数线上从 $i_0\ge1$ 出发, 每步独立地以概率 $p$ 向右走一步, 以概率 $1-p$ 向左走一步, $0\le p\le1$. 到达0时坠落. 求最终坠落的概率.
\end{exercise}
\sourceof{官方题面 \texttt{pset1.pdf} 第1页; 无公开对应官方解答, 编者独立补解}
\begin{solution}
记 $q=1-p$, $\tau_a$ 为首次到达整数 $a$ 的时刻. 所求为
\[
 \Prob_{i_0}(\tau_0<\infty)=
 \begin{cases}1,&0\le p\le\tfrac12,\\(q/p)^{i_0},&\tfrac12<p\le1.\end{cases}
\]
$p=0$ 时一直向左走, $p=1$ 时一直向右走, 两个端点直接成立. 以下设 $0<p<1$.

取整数 $N>i_0$, 先在有限区间 $[0,N]$ 中计算
$u_i^{(N)}=\Prob_i(\tau_0<\tau_N)$. 从任何内部位置都有一条不超过 $N$ 步的连续向左路径到达0, 其概率至少为 $q^N>0$. 因而每经过 $N$ 步仍未吸收的条件概率至多 $1-q^N$, 最终必在0或 $N$ 吸收. 按第一步分解,
\[
 u_0^{(N)}=1,\quad u_N^{(N)}=0,\qquad
 u_i^{(N)}=p u_{i+1}^{(N)}+q u_{i-1}^{(N)}\quad(1\le i<N).
\]
两边移项得 $p(u_{i+1}-u_i)=q(u_i-u_{i-1})$. 当 $p\ne q$ 时, 令 $s=q/p$, 相邻差为等比数列, 再代入端值得到
\[
 u_i^{(N)}=\frac{s^i-s^N}{1-s^N}.
\]
当 $p=q=1/2$ 时相邻差相等, 所以 $u_i^{(N)}=1-i/N$. 这些边界方程的解唯一: 两个解之差若有正的内部最大值, 正系数平均关系会使两个邻点取同一最大值, 沿区间传到边界而矛盾; 负的最小值同理.

事件 $\{\tau_0<\tau_N\}$ 随 $N$ 增大而递增, 其并集恰为 $\{\tau_0<\infty\}$: 一条在有限时间到达0的路径, 在此之前只经过有限多个位置. 因而所求概率是 $u_{i_0}^{(N)}$ 在 $N\to\infty$ 时的极限. $p<1/2$ 时 $s>1$, 极限为1; $p=1/2$ 时仍为1; $p>1/2$ 时 $s<1$, 极限为 $s^{i_0}$. 这里必须作有限区间极限, 不能给无穷区间另添一个未经证明的边界值.
\end{solution}

\begin{exercise}\label{mit:18212-pset1-q3}
\textbf{18.212 Spring 2019, Problem Set 1, Problem 3.}
沿用上一题的随机游走, 求从 $i_0\ge1$ 出发恰好经过 $m$ 步首次到达0的概率. 可用反射原理.
\end{exercise}
\sourceof{官方题面 \texttt{pset1.pdf} 第1页; 官方 \texttt{pset1\_solnii.pdf} 第2--3页局部 Problem 4, Sarah Wang 的解答; 本稿补全反射的起终点及偏置概率处理}
\begin{solution}
若 $m<i_0$ 或 $m-i_0$ 为奇数, 概率为0. 对其余情形, 令
$r=(m-i_0)/2$, $\ell=(m+i_0)/2$, 分别为右步、左步总数, 则
\[
 \Prob_{i_0}(\tau_0=m)
   =\frac{i_0}{m}\binom mr p^r(1-p)^\ell.
\]
首次到达0的最后一步必为从1到0的左步. 去掉它, 要数长度 $m-1$、从 $i_0$ 到1并始终在正整数中的路径. 不加正性限制时有 $\binom{m-1}{r}$ 条.

对其中曾到达0的坏路径, 找到首次到达0的时刻, 把此前的路径段关于高度0反射. 这给出从 $-i_0$ 到1的长度 $m-1$ 路径. 反向操作是把这条新路径首次到达0之前的部分再反射, 所以这是双射. 新路径的右步数为 $\ell$, 坏路径数为
$\binom{m-1}{\ell}=\binom{m-1}{r-1}$, 其中下指标负数时二项式系数为0. 因而好路径数是
\[
 \binom{m-1}{r}-\binom{m-1}{r-1}
  =\frac{m-2r}{m}\binom mr
  =\frac{i_0}{m}\binom mr.
\]
所有原来的好路径均有 $r$ 个右步及 $\ell$ 个左步, 各自概率都是 $p^r q^\ell$, 乘上路径数即得结果. 反射只用于计数; $p\ne1/2$ 时它不保持路径概率. 若 $p=0$ 或1, 用 $0^0=1$ 的空乘积约定解释公式. 特别地 $i_0=1,m=2r+1$ 时, 路径数为 Catalan 数 $\frac1{r+1}\binom{2r}{r}$.

官方文字把所用反射轴称为 $y$ 轴; 以横轴表示时间、纵轴表示位置时, 应当是高度0的横轴. 上述起终点核算给出实际使用的反射.
\end{solution}

\begin{exercise}\label{mit:18212-pset1-q10}
\textbf{18.212 Spring 2019, Problem Set 1, Problem 10.}
固定正整数 $m,n$. 证明长度为 $mn+1$ 的排列必有长度 $m+1$ 的递增子序列或长度 $n+1$ 的递减子序列.
\end{exercise}
\sourceof{官方题面 \texttt{pset1.pdf} 第2页; 官方 \texttt{pset1\_solnii.pdf} 第3页局部 Problem 6; 本稿明确各长度均以指定位置结尾}
\begin{solution}
写排列为 $w_1,\ldots,w_{mn+1}$. 令 $a_i$ 为以 $w_i$ 结尾的最长递增子序列长度, $b_i$ 为以 $w_i$ 结尾的最长递减子序列长度. 对 $i<j$, 排列元素互异, 所以或者 $w_i<w_j$, 此时在一条长度 $a_i$ 的递增子序列后接 $w_j$ 得 $a_j\ge a_i+1$; 或者 $w_i>w_j$, 此时 $b_j\ge b_i+1$. 因而所有 $(a_i,b_i)$ 两两不同.

若两种要求的长子序列都不存在, 则 $1\le a_i\le m$, $1\le b_i\le n$, 只有 $mn$ 个可能的有序对, 无法容纳 $mn+1$ 个不同的对. 矛盾. 这也是正文推论\ref{cor:es}所用的偏序长度编码.
\end{solution}

\begin{exercise}\label{mit:18212-pset1-q12}
\textbf{18.212 Spring 2019, Problem Set 1, Problem 12.}
设偏序的 Hasse 图为根树 $T$, 共 $n\ge1$ 个顶点, 根最小且祖先先于后代. 对顶点 $v$, 令 $h(v)$ 为以 $v$ 为根的子树顶点数, 包括 $v$ 自身. 证明其线性扩张数为
$n!/\prod_{v\in T}h(v)$.
\end{exercise}
\sourceof{官方题面 \texttt{pset1.pdf} 第2页; 官方 \texttt{pset1\_solnii.pdf} 第3页局部 Problem 5, Sanzeed Anwar 的解答; 本稿明确根树偏序的方向}
\begin{solution}
线性扩张是列出全部顶点的顺序, 使 $u<v$ 时 $u$ 出现在 $v$ 前面. 根必须最先出现. 假设根的各个孩子所生子树为 $T_1,\ldots,T_s$, 大小为 $n_1,\ldots,n_s$, 则 $\sum n_i=n-1$. 不同子树之间的顶点不可比, 因而剩余顺序恰由每棵子树的一种线性扩张及它们的一种交错排列确定. 保持各子树内部顺序的交错方式数为
$\binom{n-1}{n_1,\ldots,n_s}=(n-1)!/\prod_i n_i!$.

对 $n$ 归纳. $n=1$ 时两边均为1. 归纳假设给出
\[
 e(T)=\frac{(n-1)!}{\prod_i n_i!}
       \prod_i\frac{n_i!}{\prod_{v\in T_i}h(v)}
      =\frac{(n-1)!}{\prod_{v\ne\mathrm{root}}h(v)}
      =\frac{n!}{\prod_{v\in T}h(v)},
\]
最后一步用 $h(\mathrm{root})=n$. 每种交错可唯一恢复各子树内部顺序, 所以乘法计数没有重复. 三顶点根星形有2种扩张, 根链只有1种, 与公式相符. 将整个偏序反向并把每个扩张倒序, 数目不变; 这里不把树边任意定向后的所有偏序都称为同一个根树偏序.
\end{solution}

\begin{exercise}\label{mit:18212-pset1-q19}
\textbf{18.212 Spring 2019, Problem Set 1, Problem 19.}
对整数 $1\le k\le n/2$, 构造 $[n]=\{1,\ldots,n\}$ 的 $k$ 元子集到 $(n-k)$ 元子集的双射 $f$, 满足每个 $I$ 都有 $I\subseteq f(I)$.
\end{exercise}
\sourceof{官方题面 \texttt{pset1.pdf} 第4页; 无公开对应官方解答, 编者独立构造并证明逆算法}
\begin{solution}
用0--1串 $s_1\cdots s_n$ 表示子集, $s_i=1$ 表示 $i\in I$. 把1当左括号, 0当右括号, 从左向右扫描: 遇1压入栈, 遇0且栈非空时, 与栈顶的1配对并弹栈; 遇0且栈空时不配对. 记配对数为 $t$. 未配对的位依次必为
$0^a1^b$: 若一个未配对的1出现在未配对的0之前, 扫描到这个0时栈非空, 它不可能仍未配对. 因而
\[
 k=t+b,\qquad n-k=t+a,\qquad a-b=n-2k\ge0.
\]
保持全部已配对位不动, 把未配对位上的串 $0^a1^b$ 改为 $0^b1^a$, 新串所代表的集合就是 $f(I)$. 新串有 $t+a=n-k$ 个1. 因 $a\ge b$, 原来末尾的 $b$ 个未配对1仍在1的区域中, 配对1也没有改变, 所以 $I\subseteq f(I)$.

还需证明第二次扫描不会改变配对结构. 任何一对括号都不能跨过一个未配对位: 未配对0若在其内部, 在扫描该0时栈里有该对的左括号, 故这个0会被配对; 未配对1若在其内部, 外层右括号到来前必须先弹出这个1, 它也会被配对. 因而各相邻未配对位之间、以及两端的已配对位, 分成独立的平衡括号块. 每块从自身开始扫描的栈高度从不为负且最后回到0, 所有内部配对均保持. 块与块之间改后的未配对串仍为先0后1, 不会产生跨块的新配对. 所以新串的未配对位正是原来的那些位, 只是 $(a,b)$ 对调.

逆算法是对目标串执行同样的括号配对和未配对位变换, 将 $0^b1^a$ 换回 $0^a1^b$. 因而这一变换在全部0--1串上是对合, 限制到两个互补的层就得到所求双射. 例如 $n=6,k=2,I=\{1,2\}$ 的串为110000, 配对为 $(2,3),(1,4)$, 最后两位未配对, 故 $f(I)=\{1,2,5,6\}$. $n=2k$ 时 $a=b$, 双射为恒等映射.
\end{solution}

\par\medskip\noindent\textbf{Problem Set 2：匹配与分层图上的游走}

\begin{exercise}\label{mit:18212-pset2-q7}
\textbf{18.212 Spring 2019, Problem Set 2, Problem 7.}
对有限偏序集 $P$, 证明最大反链大小等于把 $P$ 分划为互不相交的链所需的最小链数. 链不必饱和.
\end{exercise}
\sourceof{官方题面 \texttt{pset2.pdf} 第2页; 无公开对应官方解答, 编者独立补解; 使用正文定理\ref{thm:konig}的完整证明}
\begin{solution}
令 $N=|P|$. $N=0$ 时两个数均为0. 以下设 $N>0$. 建立二部图: 每个元素 $x$ 在左右各有副本 $x_L,x_R$, 当且仅当 $x<y$ 时连边 $x_Ly_R$. 记最大匹配大小为 $\nu$.

把匹配边 $x_Ly_R$ 解释成 $P$ 中的有向边 $x\to y$. 每个点入度和出度至多1, 又因严格偏序不能形成有向圈, 这些边把所有元素分成互不相交的有向路径和孤立点. 每条路径的元素由传递性构成一条链. $N$ 个点加上 $\nu$ 条这样的边, 所得链数为 $N-\nu$. 反过来, 一个 $c$ 链分划, 在每条链按递增顺序连接相邻元素, 就给出二部图中大小 $N-c$ 的匹配: 左副本不重复作为前项, 右副本不重复作为后项. 因而最小链数恰为 $N-\nu$.

由 Kőnig 定理\ref{thm:konig}, 二部图有大小为 $\nu$ 的顶点覆盖 $C$. 取
$A=\{x\in P:x_L\notin C,\ x_R\notin C\}$. $C$ 的每个副本最多排除一个原元素, 故 $|A|\ge N-\nu$. 若 $x,y\in A$ 且 $x<y$, 边 $x_Ly_R$ 的两个端点都不在覆盖中, 矛盾; 所以 $A$ 是反链. 每条链至多与一条反链相交于一个元素, 因而任何反链大小又不超过任何链分划的链数 $N-\nu$. 两端合并即得等式, 也构造了达到等号的链分划和反链. 正文定理\ref{thm:dilworth}与本题对应.
\end{solution}

\begin{exercise}\label{mit:18212-pset2-q9}
\textbf{18.212 Spring 2019, Problem Set 2, Problem 9.}
Young 格由所有整数分拆的 Young 图按包含关系排序而成, 相邻图恰差一个方格. 令 $W_n$ 为其 Hasse 图上从空图出发、经过 $2n$ 步又回到空图的游走数, 每步可向上加格或向下删格. 证明 $W_n$ 等于完全图 $K_{2n}$ 的完美匹配数, 并给出闭式. 例如 $W_2=3$.
\end{exercise}
\sourceof{官方题面 \texttt{pset2.pdf} 第2页; 官方 \texttt{pset2\_solni.pdf} 第2页局部 Problem 4, Yogeshwar Velingker 的解答; 本稿补证算子关系及其计数转换}
\begin{solution}
Young 图可视为格点 $(i,j)\in\N^2$ 的有限向下封闭集合: 若 $(i,j)$ 在图中, 则左方和上方的格子也在图中. 给每个图 $\lambda$ 一个形式基向量 $v_\lambda$. 定义 $U v_\lambda$ 为加一个格后所有图的基向量之和, $D v_\lambda$ 为删一个格后所有图之和. 那么 $(U+D)^r v_\varnothing$ 中 $v_\mu$ 的系数, 正是长度 $r$、从空图到 $\mu$ 的游走数.

证明
\begin{equation}\label{eq:18212-young-commutator}
 DU-UD=I.
\end{equation}
对两个不同但同样大小的图 $\lambda,\mu$, $DU$ 的相应系数是共同上覆盖的数目, $UD$ 的系数是共同下覆盖的数目. 若它们恰各有一个格不属于对方, 则共同上覆盖唯一为 $\lambda\cup\mu$, 共同下覆盖唯一为 $\lambda\cap\mu$, 两者仍是向下封闭集合; 若差格更多, 两种覆盖都不存在. 因而非对角系数相等. 对角系数之差是可加角的数目减可删角的数目. 若非空分拆的正行长有 $s$ 个不同值, 每个相同行长块的最后一行末端给出一个可删角, 共 $s$ 个. 第一行可在末端加格, 以后每次行长下降可在新块的第一行末端加格, 还可在最底端另加一行; 共 $s+1$ 个. 空图也有1个可加角而无可删角. 这就证明了式\eqref{eq:18212-young-commutator}.

为提取回到空图的系数, 仅使用形式幂级数. 每个 $t^r$ 系数只涉及有限多个大小不超过 $r$ 的图, 所以不需收敛假设. 由上述交换关系归纳得
$DU^k=U^kD+kU^{k-1}$, 从而
$D e^{tU}=e^{tU}(D+tI)$. 现在有正规排序公式
\[
 e^{t(U+D)}=e^{t^2/2}e^{tU}e^{tD}.
\]
验证如下: 右边记为 $F(t)$, 用刚才的交换关系求形式导数得
$F'(t)=(U+D)F(t)$, 且 $F(0)=I$. 若 $F=\sum F_r t^r$, 方程给出 $(r+1)F_{r+1}=(U+D)F_r$, 唯一确定 $F_r=(U+D)^r/r!$, 所以等式成立.

因 $Dv_\varnothing=0$, 右边作用于空图时 $e^{tD}$ 消失; $U$ 严格增加大小, 所以 $e^{tU}v_\varnothing$ 的空图系数只有常数项1. 因而
\[
 \sum_{r\ge0}\frac{t^r}{r!}
     [v_\varnothing](U+D)^r v_\varnothing
       =e^{t^2/2},\qquad
 W_n=\frac{(2n)!}{2^n n!}=(2n-1)!!.
\]
另一方面, 给 $K_{2n}$ 的顶点1选择匹配伙伴有 $2n-1$ 种, 剩余顶点递归匹配, 得 $(2n-1)(2n-3)\cdots1$ 种, 恰与 $W_n$ 相同. $n=0$ 时空游走与空匹配各1种. $n=2$ 时三条游走的中间两步分别经过图 $(2)$、图 $(1,1)$、空图, 对应题面所列的3条; 此处证明的是数目等式, 不把算子等式宣称为已构造了一一对应.
\end{solution}

\par\medskip\noindent\textbf{Problem Set 3：树计数、有向圈、电网络与停车函数}

\begin{exercise}\label{mit:18212-pset3-q1}
\textbf{18.212 Spring 2019, Problem Set 3, Problem 1.}
给出完全二部图 $K_{m,n}$ 生成树数目 $m^{n-1}n^{m-1}$ 的双射证明, 其中 $m,n\ge1$.
\end{exercise}
\sourceof{官方题面 \texttt{pset3.pdf} 第1页; 官方 \texttt{pset3\_soln.pdf} 第1页局部 Problem 1, Sanzeed Anwar 的解答; 本稿补全删除规则与任意词对的逆编码}
\begin{solution}
两部为分别带全序的集合 $A,B$, 大小为 $m,n$. 要把生成树编码成一对词
$(\alpha,\beta)\in A^{n-1}\times B^{m-1}$; 空词长度为0.

给定树, 设尚未删除的两部大小为 $a,b$. 当 $a+b>2$ 时, 若 $a\ge b$ 就从 $A$ 删除标号最小的叶, 将其邻点记到词 $\beta$ 末尾; 若 $b>a$ 就从 $B$ 删除最小的叶, 将邻点记到 $\alpha$ 末尾. 这条规则总能执行: 当前图仍为二部树, 无孤立点, 且有 $a+b-1$ 条边. 若 $a\ge b$ 而 $A$ 中没有叶, 则所有 $A$ 顶点度数至少2, 给出 $a+b-1\ge2a$, 与 $b\le a$ 矛盾; $b>a$ 时同理. 最后只剩两部各一个点及其连接边. 一共删去 $m-1$ 个 $A$ 点、$n-1$ 个 $B$ 点, 所得词长正是规定长度. 每次删哪一部只由当前 $a,b$ 决定, 所以删部的顺序不依赖具体树.

逆算法从任意词对开始. 沿同一删部顺序, 若本次删 $A$, 在尚存的 $A$ 点中取没有出现在 $\alpha$ 未读部分的最小标号 $u$, 把它连向 $\beta$ 的下一个未读字母, 删除 $u$ 并读掉这个 $\beta$ 字母. 若本次删 $B$, 在尚存的 $B$ 点中取没有出现在 $\beta$ 未读部分的最小点, 连向 $\alpha$ 的下一个字母并作相应删除. 当各部仅剩一个点时连上最后的边.

需要核对候选点及邻点都合法. 当前 $\alpha$ 未读长度为 $b-1$, $\beta$ 未读长度为 $a-1$. 若 $a\ge b$, $A$ 的 $a$ 个点不可能全出现在长度 $b-1<a$ 的词中, 故所取 $A$ 点存在; 若 $b>a$, 另一个判断同理. 任何点只在其标号不再出现在将来词中时才被删除, 所以被读出的邻点标号仍存活. 若 $a\ge b$ 且尚未结束, 则 $a\ge2$, 有下一个 $\beta$ 字母可读, 对另一部同理.

所得图是树: 每个删掉的点连向一个之后才删除的点, 最后两个点相连, 顺着删除时间前进必到达这条末边, 故连通且有 $m+n-1$ 条边. 在算法的任何中间阶段, 一个剩余点在剩余树中的度数等于1加上它在相应未读词中的出现次数. 因而候选点恰是指定部的叶, 最小标号规则也完全一致; 重新编码会读回原词对. 反之, 一棵树的编码中, 一个点在未读邻点记录中的出现次数为当前度数减1, 其最后一条边不在它删除前记录为它自身; 解码因此逐步恢复原来的叶及边. 两个算法互逆. 词对数为 $m^{n-1}n^{m-1}$, 即所求. $m=1$ 或 $n=1$ 时一部分词为空, 另一部分只有唯一字母可用, 正确给出唯一的星形树.
\end{solution}

\begin{exercise}\label{mit:18212-pset3-q2}
\textbf{18.212 Spring 2019, Problem Set 3, Problem 2.}
求完全三部图 $K_{m,n,k}$ 的生成树数, 并证明. 三部大小为正整数 $m,n,k$, 两点当且仅当分属不同部时相邻.
\end{exercise}
\sourceof{官方题面 \texttt{pset3.pdf} 第1页; 无公开对应官方解答, 编者独立补解}
\begin{solution}
令 $N=m+n+k$. 答案为
\[
 \tau(K_{m,n,k})
  =N(n+k)^{m-1}(m+k)^{n-1}(m+n)^{k-1}.
\]
先从矩阵树定理\ref{thm:matrix-tree}推出稍后也会使用的谱计数式. 对 $N$ 点连通图, 拉普拉斯矩阵 $L$ 的特征值为 $0,\lambda_2,\ldots,\lambda_N$, 其中 $\lambda_i>0$: 二次型为各边差平方之和, 零空间在连通图上恰为常数向量. 在 $\det(L+tI)$ 中取 $t$ 的一次项, 按选择对角项的行列式展开, 得所有主余子式的和 $N\tau(G)$. 由特征值乘积, 同一系数又是 $\prod_{i=2}^N\lambda_i$. 所以
\begin{equation}\label{eq:18212-spectral-tree}
 \tau(G)=\frac1N\prod_{i=2}^N\lambda_i.
\end{equation}

现在计算三部图的谱. 只在第一部取非零值且该部坐标和为0的向量, 经 $L$ 作用后乘以 $N-m=n+k$, 维数为 $m-1$; 另两部类似, 得特征值 $m+k$、$m+n$, 重数分别为 $n-1,k-1$. 剩余三维空间由每一部上的常值向量组成. 若三部常值为 $x_1,x_2,x_3$, 记 $S=mx_1+nx_2+kx_3$, 则
\[
 (Lx)_i=Nx_i-S\quad(i=1,2,3).
\]
全常值向量的特征值为0; $S=0$ 的二维子空间上特征值为 $N$. 所列空间维数和为 $N$, 覆盖全部特征值. 代入式\eqref{eq:18212-spectral-tree}, 得 $N^2/N$ 乘上三部内部的特征值乘积, 即上述公式. 例如 $m=n=k=1$ 给出 $\tau(K_3)=3$.
\end{solution}

\begin{exercise}\label{mit:18212-pset3-q3}
\textbf{18.212 Spring 2019, Problem Set 3, Problem 3.}
设 $G$ 为 $n\ge1$ 个顶点上的简单图, $\bar G$ 为同一顶点集上的补图. 给 $G$ 加一个新点0并把它连向所有旧点, 得 $G^+$. 定义
$F_G(x)=\sum_T x^{\deg_T(0)-1}$, 和遍历 $G^+$ 的生成树. 证明
$F_{\bar G}(x)=(-1)^{n-1}F_G(-x-n)$.
\end{exercise}
\sourceof{官方题面 \texttt{pset3.pdf} 第1页; 无公开对应官方解答, 编者独立补解}
\begin{solution}
把每条0到旧点的边赋权 $x$, 旧图的边赋权1. 对 $x>0$, 加权矩阵树定理\ref{thm:matrix-tree}给出
\[
 xF_G(x)=\sum_T x^{\deg_T(0)}=\det(xI_n+L_G),
\]
因为删除0的行列后, 各个旧顶点的对角元加上 $x$, 其余仍是 $L_G$. 两边都是多项式, 在无穷多个正实数处相等, 因而此等式为多项式恒等式.

$L_G$ 为实对称矩阵且消去常数向量, 可选一组正交特征向量, 第一根为常数向量, 其余在常数向量的正交补上, 相应特征值记为 $\lambda_2,\ldots,\lambda_n$. 即使 $G$ 不连通, 也允许某些 $\lambda_i=0$. 因 $n\ge1$, 每棵 $G^+$ 的生成树必须包含与0相接的边, 所以 $F_G$ 定义中的指数非负. 在多项式环中约去上述等式的因子 $x$, 得
\[
 F_G(x)=\prod_{i=2}^n(x+\lambda_i).
\]
补图满足 $L_{\bar G}=nI-J-L_G$, 其中 $J$ 是全1矩阵: 度数与补度数之和为 $n-1$, 非对角邻接互补. 在常数向量的正交补上 $J=0$, 因而补图对应特征值为 $n-\lambda_i$. 于是
\[
 F_{\bar G}(x)=\prod_{i=2}^n(x+n-\lambda_i)
   =(-1)^{n-1}\prod_{i=2}^n(-x-n+\lambda_i),
\]
即所求. $n=1$ 时两个乘积都是空乘积1.
\end{solution}

\begin{exercise}\label{mit:18212-pset3-q4}
\textbf{18.212 Spring 2019, Problem Set 3, Problem 4.}
证明以下两条恒等式, $n\ge0$:
\begin{align*}
 (x+y)^n&=\sum_{k=0}^n\binom nk y(y+kz)^{k-1}(x-kz)^{n-k},\\
 (x+y)(x+y+nz)^{n-1}
   &=\sum_{k=0}^n\binom nk x(x+kz)^{k-1}
                       y(y+(n-k)z)^{n-k-1}.
\end{align*}
在 $k=0$ 或 $n-k=0$ 的端点, 相应因子 $t(t+0z)^{-1}$ 按约去后的多项式1解释; 第二式 $n=0$ 的左边也解释为1. 原题提示可使用有向矩阵树定理.
\end{exercise}
\sourceof{官方题面 \texttt{pset3.pdf} 第1页的两条恒等式; 官方 \texttt{pset3\_soln.pdf} 第2--3页局部 Problem 3, Congyue Deng 与 Ganatra 的解答; 本稿补全权重、方向、端点及第一式的独立代数证明}
\begin{solution}
记 $H_0(t,z)=1$, $H_k(t,z)=t(t+kz)^{k-1}$ $(k\ge1)$. 这是统一处理端点的多项式记号.

先证第二式. 建立有向图, 顶点为 $A,B$ 及 $n$ 个普通点. 弧 $A\to B$ 权为1, $A$ 到每个普通点的弧权为 $x$, $B$ 到每个普通点的弧权为 $y$, 普通点之间双向弧权为 $z$, 无其他弧. 暂设 $x,y,z>0$. 数根为 $A$、从根向外的有向生成树; 用定理\ref{thm:directedmatrix-tree}对反向图的出度拉普拉斯取删根余子式, 等价于原图按入弧权求和的拉普拉斯余子式. 按顶点 $B$、普通点排序, 矩阵形如
\[
 \begin{pmatrix}1&0\\-y\mathbf1&(x+y+nz)I_n-zJ_n\end{pmatrix}.
\]
当 $n\ge1$, 右下块在常数方向的特征值为 $x+y$, 在其正交补上的特征值为 $x+y+nz$, 故总权为 $H_n(x+y,z)$. $n=0$ 时只有弧 $A\to B$, 总权为1.

每棵这样的树必含 $A\to B$, 因为 $B$ 的唯一入弧是它. 删去此弧后, 得到以 $A$ 和 $B$ 为根的两个分支. 若 $A$ 分支恰含 $k$ 个普通点, 有 $\binom nk$ 种选择. 一个根连向 $k$ 个普通点、根弧权为 $t$、普通点间双向弧权为 $z$ 的出树总权是 $H_k(t,z)$: 当 $k\ge1$ 时相应删根矩阵为 $(t+kz)I_k-zJ_k$, 特征值为 $t$ 及 $k-1$ 个 $t+kz$; $k=0$ 时空树权为1. 因而这两个分支权乘积为 $H_k(x,z)H_{n-k}(y,z)$. 删边及加回弧 $A\to B$ 互逆, 得到
\[
 H_n(x+y,z)=\sum_{k=0}^n\binom nk H_k(x,z)H_{n-k}(y,z).
\]
两边是三变量多项式, 在全部正实数参数上相等, 遂为多项式恒等式, 得证第二式.

第一式采用能直接核对端点的归纳法. 设
$F_n(x,y,z)=\sum_{k=0}^n\binom nk H_k(y,z)(x-kz)^{n-k}$. $F_0=1$. 对 $x$ 求导并用 $(n-k)\binom nk=n\binom{n-1}k$, 有
\[
 \frac{\partial F_n}{\partial x}=nF_{n-1}(x,y,z).
\]
若 $F_{n-1}=(x+y)^{n-1}$, 则 $F_n-(x+y)^n$ 与 $x$ 无关. 对 $n\ge1$, 代入 $x=-y$ 得
\[
 F_n(-y,y,z)=y\sum_{k=0}^n(-1)^{n-k}\binom nk(y+kz)^{n-1}=0.
\]
这里包括 $k=0$ 项: 它在左边是 $(-y)^n$, 在右边是 $y(-1)^n y^{n-1}$. 等式最后一步来自有限差分: 若 $P(k)$ 是次数小于 $n$ 的多项式, $\Delta P(k)=P(k+1)-P(k)$ 每次使次数至少降低1, 故 $\Delta^nP(0)=0$; 反复展开差分给出 $\Delta^nP(0)=\sum_{k=0}^n(-1)^{n-k}\binom nkP(k)$. 取 $P(k)=(y+kz)^{n-1}$ 即可. 同时 $(-y+y)^n=0$, 所以与 $x$ 无关的差为0, 完成归纳及第一式的证明.
\end{solution}

\begin{exercise}\label{mit:18212-pset3-q5}
\textbf{18.212 Spring 2019, Problem Set 3, Problem 5.}
$n\ge1$ 维立方体图 $Q_n$ 的顶点是 $\{0,1\}^n$, 两点恰在一个坐标不同则相邻. 把每条边换成一对反向弧, 求经过每条弧恰好一次的有向 Euler 圈数.
\end{exercise}
\sourceof{官方题面 \texttt{pset3.pdf} 第2页; 官方 \texttt{pset3\_soln.pdf} 第3页局部 Problem 4, Sophia Xia 的解答; 本稿补证立方体树数并明确圈的计数惯例}
\begin{solution}
这里把圈视为弧的循环序列, 仅识别循环转动, 不识别反向或图的自同构. 等价地, 预先指定一条弧 $e_0$ 并将它固定为第一条弧, 计线性排列. 每条弧恰出现一次, 因而每个循环序列恰有一个以 $e_0$ 起首的代表. 在这一惯例下, 答案为
\[
 2^{\,2^n-n-1}
 \left(\prod_{k=1}^n k^{\binom nk}\right)
 ((n-1)!)^{2^n}.
\]

计算生成树数. 对每个坐标子集 $S\subseteq[n]$, 定义字符函数
$\chi_S(v)=(-1)^{\sum_{i\in S}v_i}$. 若 $S\ne T$, 在 $S\mathbin\triangle T$ 中取一个坐标并配对该坐标翻转的两个顶点, 则 $\sum_v\chi_S(v)\chi_T(v)=0$; 每个字符平方和为 $2^n$. 这些 $2^n$ 个非零正交函数因此构成全部顶点函数空间的基. 翻转不在 $S$ 的坐标时值不变, 翻转在 $S$ 的坐标时值反号, 所以
\[
 L\chi_S=2|S|\chi_S.
\]
特征值 $2k$ 的重数为 $\binom nk$, $k=0,\ldots,n$. 图连通, 由式\eqref{eq:18212-spectral-tree}得
\[
 \tau(Q_n)=\frac1{2^n}\prod_{k=1}^n(2k)^{\binom nk}
      =2^{\,2^n-n-1}\prod_{k=1}^n k^{\binom nk},
\]
因为 $\sum_{k=1}^n\binom nk=2^n-1$.

双向化后每点入度、出度都是 $n$, 且图强连通. 对 $e_0$ 的起点 $r$, 无向生成树朝 $r$ 定向与根为 $r$ 的入树一一对应, 故入树数也是 $\tau(Q_n)$. 由 BEST 定理\ref{thm:best}, 固定首弧的 Euler 圈数为
$\tau(Q_n)\prod_v(d^+(v)-1)!=\tau(Q_n)((n-1)!)^{2^n}$, 得所求. 若固定起点而允许任意首弧, 应再乘 $n$; 若每个任意起点的线性弧序列都单独计数, 应乘弧总数 $n2^n$. $n=1$ 时一个双向边有1个循环序列, $n=2$ 时结果为4, 与该惯例一致.
\end{solution}

\begin{exercise}\label{mit:18212-pset3-q6}
\textbf{18.212 Spring 2019, Problem Set 3, Problem 6.}
$n\ge1$ 维立方体图的每条边电阻均为 $1\,\Omega$. 求一对相反顶点之间的有效电阻.
\end{exercise}
\sourceof{官方题面 \texttt{pset3.pdf} 第2页; 无公开对应官方解答, 编者独立补解}
\begin{solution}
端点取 $a=(0,\ldots,0)$, $b=(1,\ldots,1)$. 电导均为1. 注入 $a$ 一单位电流、从 $b$ 流出一单位, 取 $\phi(b)=0$. 图有限、连通且 $a\ne b$, 接地电位由正文命题\ref{prop:potential}唯一确定. 任意坐标置换保持端点及网络不变, 唯一性迫使电位只依赖 Hamming 重量 $k$, 记为 $\phi_k$. 这也说明把同层点合并不产生虚假的同层电流.

第 $k$ 层到第 $k+1$ 层的边数为
$\binom nk(n-k)=n\binom{n-1}k$ $(0\le k<n)$. 在重量不超过 $k$ 的全部顶点上求和 Kirchhoff 方程, 内部边的电流成对抵消, 只剩穿过这一层边界的电流, 总量为1. 所有这些边的电流均为 $\phi_k-\phi_{k+1}$, 故
\[
 \phi_k-\phi_{k+1}=\frac1{n\binom{n-1}k}.
\]
逐层相加, 一单位电流下的端点电压差就是有效电阻:
\[
 R_{\mathrm{eff}}(a,b)=\phi_0-\phi_n
       =\frac1n\sum_{k=0}^{n-1}\frac1{\binom{n-1}k}\ \Omega.
\]
$n=1,2,3$ 分别得到 $1,1,5/6\,\Omega$. 若把此网络看作简单随机游走, 加权度数总和为 $\mathcal V=n2^n$, 通勤公式\ref{thm:commute}给出
$\E_a\tau_b+\E_b\tau_a=2^n\sum_{k=0}^{n-1}\binom{n-1}k^{-1}$; 取坐标全部翻转的对称性, 两个单向期望相等. 这一核对使用有限连通网络和不同的端点, 不把端点作为游走中每次都被吸收的两个边界来计算通勤时间.
\end{solution}

\begin{exercise}\label{mit:18212-pset3-q7}
\textbf{18.212 Spring 2019, Problem Set 3, Problem 7.}
$n\ge1$ 张卡片编号为 $1,\ldots,n$. Alice 最初有 $k$ 张, Ben 有 $n-k$ 张, $0\le k\le n$. 每轮独立均匀选一个编号, 持有该卡的人把它交给另一人. 一人集齐全部卡片时游戏结束. 求 Alice 获胜的概率.
\end{exercise}
\sourceof{官方题面 \texttt{pset3.pdf} 第2页; 无公开对应官方解答, 编者独立补解}
\begin{solution}
用 Alice 所持卡片的指示向量表示状态, 每轮恰翻转一个均匀选择的坐标, 所以这是 $Q_n$ 上的简单随机游走, 空集和全集分别为失败和成功边界. 因卡片编号的置换对称性, 胜率只依赖持有数量 $k$, 记为 $h_k$.

游戏几乎必结束: 从任何非终止状态, 把 Alice 尚未持有的卡片各转交一次, 至多 $n$ 步可到全集, 特定的这些选择概率至少为 $n^{-n}>0$; 若更早到达终止状态也已结束. 分成长度 $n$ 的时间段, 未结束概率至多 $(1-n^{-n})^r$ 而趋于0. 因而第一步分解的解确为获胜概率, 满足
\[
 h_0=0,\quad h_n=1,\qquad
 h_k=\frac{k}{n}h_{k-1}+\frac{n-k}{n}h_{k+1}\quad(1\le k<n).
\]
令 $\Delta_i=h_i-h_{i-1}$. 方程给出
$\Delta_{k+1}=\frac{k}{n-k}\Delta_k$, 从而
$\Delta_i=\Delta_1/\binom{n-1}{i-1}$. 用 $\sum_{i=1}^n\Delta_i=1$ 求得
\[
 h_k=\frac{\displaystyle\sum_{j=0}^{k-1}\binom{n-1}j^{-1}}
            {\displaystyle\sum_{j=0}^{n-1}\binom{n-1}j^{-1}}.
\]
$k=0$ 的分子是空和0, $k=n$ 时等于1; $n=1$ 时没有内部方程, 公式仍适用. $n=3,k=1$ 时胜率为 $2/5$, 并非 $k/n$. 对称性给出 $h_k+h_{n-k}=1$, 也可由二项式系数的对称性直接核对.
\end{solution}

\begin{exercise}\label{mit:18212-pset3-q8}
\textbf{18.212 Spring 2019, Problem Set 3, Problem 8.}
道路上有编号 $1,\ldots,n$ 的 $n$ 个停车位, 汽车按编号 $1,\ldots,n$ 到达. 第 $i$ 辆车希望停在 $f(i)$, 实际选择尚空且编号不小于 $f(i)$ 的最小车位; 没有这样的车位就失败. $f:[n]\to[n]$ 称为停车函数, 当所有车都能停车. 证明以下三条件等价:
\begin{enumerate}[label=(\arabic*)]
\item $f$ 是停车函数;
\item 对 $k=1,\ldots,n$, $\#\{i:f(i)\ge n-k+1\}\le k$;
\item 存在排列 $w\in S_n$, 使每个 $i$ 有 $f(i)\le w(i)$.
\end{enumerate}
\end{exercise}
\sourceof{官方题面 \texttt{pset3.pdf} 第2页; 官方 \texttt{pset3\_soln.pdf} 第1页局部 Problem 2, Fadi Atieh 的解答; 官方解答误写 ``at least $k$'', 本稿按题面改为至多 $k$ 并补全贪心交换论证}
\begin{solution}
(1)推出(3): 令 $w(i)$ 为第 $i$ 辆车实际占用的车位, 所有车位恰好用一次, 且 $w(i)\ge f(i)$.

(3)推出(2): 希望车位在最后 $k$ 个位置的车, 在 $w$ 中都必须分配到这最后 $k$ 个位置, 所以至多有 $k$ 辆.

(2)推出(3): 将车按希望车位非递减排列, 得
$a_1\le\cdots\le a_n$, 相等时按车号排以保证确定性. 必有 $a_j\le j$: 否则 $a_j\ge j+1$, 有至少 $n-j+1$ 辆车希望在 $j+1,\ldots,n$ 中, 超出最后 $n-j$ 个车位的容量, 与(2)矛盾. 把这一排序中的第 $j$ 辆车分配给车位 $j$, 就得到满足 $f(i)\le w(i)$ 的排列.

最后证(3)推出(1), 这一步要保证题设原来的到达顺序也能成功. 在处理下一辆车 $i$ 之前, 保持一个把未到车辆双射分配到尚空车位的可行分配, 每车分配位不小于其希望位. 起初由 $w$ 给出. 设它把 $i$ 分给 $s$, 贪心停车选的最小空位为 $t\ge f(i)$; 因 $s$ 是可行空位, $t$ 一定存在且 $t\le s$. 若 $t=s$, 直接删除车辆及车位. 若 $t<s$, 则在现有可行分配中有另一辆车 $j$ 被分给 $t$, 因而 $f(j)\le t\le s$. 将 $j$ 改分给 $s$, 将 $i$ 分给 $t$ 再删除 $i,t$, 其余可行性不变. 归纳下去所有车均成功. 三条件于是等价.

另一个随后使用的写法是
$\#\{i:f(i)\le k\}\ge k$ $(1\le k\le n)$: 对 $k<n$, 它是(2)中最后 $n-k$ 个车位条件的补集形式; $k=n$ 自动成立. 官方解答的 ``at least'' 与原题不一致, 例如所有 $f(i)=1$ 在 $n>1$ 时确为停车函数, 却不满足最后一个车位至少有一辆车希望停入的错误要求.
\end{solution}

\begin{exercise}\label{mit:18212-pset3-q9}
\textbf{18.212 Spring 2019, Problem Set 3, Problem 9.}
构造以下四个集合之间的双射, $n\ge1$:
\begin{enumerate}[label=(\arabic*)]
\item 顶点标号为 $0,1,\ldots,n$ 的树, 以0作为指定根;
\item $n$ 个顶点的平面二叉树, 每点带一个不同的 $1,\ldots,n$ 标号, 每个左孩子的标号大于其父亲, 对右孩子不作此限制;
\item 长度 $2n$ 的 Dyck 路, 即有 $n$ 个上步 $U$、$n$ 个下步 $D$, 每个前缀中上步不少于下步; 上步分别带标号 $1,\ldots,n$, 每两个相邻上步的后一标号大于前一标号;
\item 大小为 $n$ 的停车函数.
\end{enumerate}
\end{exercise}
\sourceof{官方题面 \texttt{pset3.pdf} 第2页; 官方 \texttt{pset3\_soln.pdf} 第4页局部 Problem 6, Vanshika 的解答只详述二叉树与 Dyck 路, 对停车函数仅作提示; 本稿独立补全其余可逆算法}
\begin{solution}
给出一条双射链 $(1)\leftrightarrow(4)\leftrightarrow(3)\leftrightarrow(2)$. 为递归方便, 也约定 $n=0$ 时只有单独的根、空停车函数、空路和空二叉树.

\textbf{停车函数与带标号 Dyck 路.}
对 $k=1,\ldots,n$, 把满足 $f(i)=k$ 的车号按递增顺序写成一串带标号的上步, 然后写一个下步. 全部连接起来. 在第 $k$ 个下步后, 高度为 $\#\{i:f(i)\le k\}-k\ge0$, 由习题\ref{mit:18212-pset3-q8}的等价条件保证. 每串上步内部标号递增, 各串由下步隔开, 因而得到(3). 逆向时, 对标号为 $i$ 的上步, 令 $f(i)$ 等于从该上步起向右遇到的第一个下步的序号; 等价地, 它位于第 $f(i)$ 个下步之前、前一个下步之后的上步块中. 每个上步后总有一个下步, 所以 $f(i)\in[n]$. Dyck 前缀条件给出 $\#\{i:f(i)\le k\}\ge k$, 因而 $f$ 是停车函数. 块内递增条件使逆向读出再正向排序不会改变上步次序, 所以两算法互逆.

\textbf{平面二叉树与带标号 Dyck 路.}
空树对应空字. 根标号为 $v$、左子树为 $L$、右子树为 $R$ 时, 递归写
$W(T)=U_vW(L)DW(R)$. 这是一条 Dyck 路. 如果左子树非空, 开头两个上步的标号分别为 $v$ 和左孩子, 恰满足递增条件; 其他相邻上步或者同在 $W(L)$ 中, 或者同在 $W(R)$ 中, 由归纳满足条件. 逆向时, 非空 Dyck 字有唯一首次返回高度0的分解 $U_v A D B$, 其中 $A,B$ 都是带标号 Dyck 字. 递归解码为根 $v$、左子树解码 $A$、右子树解码 $B$. 若 $A$ 非空, 它的首标号大于 $v$, 所以左孩子条件成立. 两个递归分解逐步互逆.

\textbf{根树与停车函数.}
以下构造还为下一题保存反序统计. 对任意有限标号集作重编号时, 都明确说明根和其余点的顺序; 不假定汽车标号必须就是同一个子树的顶点标号.

给定根为0的树 $T$ 的顶点集为 $\{0,1,\ldots,n\}$, 找到根的分支中包含最大标号 $n$ 的那一支. 记其顶点集为 $B$, 根邻点为 $v$, $j=|B|$, 并令
\[
 r=\#\{b\in B:b<v\},\qquad C=[n]\setminus B.
\]
$v$ 是 $B$ 中第 $r+1$ 小的标号, 故 $0\le r<j$. 把 $T[B]$ 以 $v$ 为根, 重编号 $v\mapsto0$, 再将 $B\setminus\{v\}$ 按递增顺序编号为 $1,\ldots,j-1$, 得内树 $T_L$. 把 $T[\{0\}\cup C]$ 保持根0, 将 $C$ 递增编号为 $1,\ldots,n-j$, 得外树 $T_R$. 递归得到停车函数 $g$ 与 $h$, 大小分别为 $j-1,n-j$.

令 $S=B\setminus\{n\}$, 按递增顺序写 $S=\{s_1<\cdots<s_{j-1}\}$, $C=\{c_1<\cdots<c_{n-j}\}$. 定义
\begin{equation}\label{eq:18212-tree-parking-map}
 f(s_\ell)=g(\ell),\qquad f(c_\ell)=j+h(\ell),\qquad f(n)=j-r.
\end{equation}
这里内树的非根顶点取 $B\setminus\{v\}$, 而接收内停车函数的车号取 $B\setminus\{n\}$; 两者都以递增的序位对应 $[j-1]$, 这一差别是构造的一部分. 车 $S$ 只需用 $1,\ldots,j-1$ 的位置, 车 $C$ 的希望位大于 $j$, 只需用 $j+1,\ldots,n$ 的位置. 两群车虽然按原车号交错到达, 互不占用对方的区域, 在各自区域内部的到达顺序恰是递归停车函数的顺序. 因而前 $n-1$ 辆车停车后, 唯一空位是 $j$. 最后一辆车 $n$ 希望位为 $j-r\le j$, 正好停入 $j$, $f$ 是停车函数.

现在给出逆算法. 给定 $f$, 让前 $n-1$ 辆车按原顺序停车, 记唯一空位为 $j$. 因全部车能停, 必有 $1\le f(n)\le j$, 令 $r=j-f(n)$. 设 $S$ 是停在 $j$ 左侧的车号集合, $C$ 是停在右侧的车号集合, 则 $|S|=j-1$, $|C|=n-j$. 对任意停在右侧的车, 必有 $f(c)>j$: $j$ 始终空着, 若 $f(c)\le j$, 到达时就会停在不大于 $j$ 的空位, 矛盾. 停在左侧的车则有 $f(s)\le j-1$. 因而按递增车号读取 $f|_S$ 得停车函数 $g$ 定义于 $[j-1]$, 按递增车号读取 $f|_C-j$ 得停车函数 $h$ 定义于 $[n-j]$. 限制到各区域后, 已发生的停车过程就是这两个小函数的成功过程.

令 $B=S\cup\{n\}$, 取 $B$ 中第 $r+1$ 小的点为 $v$. 递归反解 $g$ 得到根为0的内树, 将其根改标为 $v$, 非根点按序改标为 $B\setminus\{v\}$. 递归反解 $h$, 保持根0, 非根点按序改标为 $C$. 合并两树并加边 $0v$, 就得到 $T$. 它是树, 因两部分顶点集不相交, 加入 $0v$ 恰将两树连接.

对大小归纳证明互逆. 从树出发, 正算法使前 $n-1$ 辆车唯一留空 $j$, 左右车号正是 $S,C$, 最后车的希望位给回 $r$. 逆算法恢复相同 $B,v$, 由归纳恢复两个子树. 从停车函数出发, 逆算法产生的分支 $B$ 包含 $n$, 所以正算法选回它, 又由归纳取得相同 $g,h$, 式\eqref{eq:18212-tree-parking-map}恢复每个 $f(i)$. 空对象为归纳基点, 故这是真正可执行的双射, 四个集合间的双射链完成.

例如 $n=2$ 时, 三棵树 $\{01,02\}$、$\{01,12\}$、$\{02,12\}$ 分别对应 $(2,1)$、$(1,2)$、$(1,1)$, 这里 $01$ 表示边 $\{0,1\}$. $n=3$ 时树 $\{03,13,23\}$ 的 $B=\{1,2,3\}$、$v=3$、$j=3,r=2$, 内树是两叶根星, 对应 $g=(2,1)$, 最终得到 $f=(2,1,1)$. 这两个例子也用于下面的统计核对.
\end{solution}

\begin{exercise}\label{mit:18212-pset3-q10}
\textbf{18.212 Spring 2019, Problem Set 3, Problem 10.}
对 $\{0,1,\ldots,n\}$ 上以0为根的树 $T$, 反序对是满足 $1\le i<j\le n$ 且 $j$ 为 $i$ 的真祖先的有序对 $(i,j)$. 记其数为 $\operatorname{inv}(T)$, 树反序多项式为
$I_n(x)=\sum_T x^{\operatorname{inv}(T)}$. 证明
\[
 I_n(x)=\sum_{f\text{ 为停车函数}}
             x^{\binom{n+1}2-\sum_{i=1}^n f(i)}.
\]
\end{exercise}
\sourceof{官方题面 \texttt{pset3.pdf} 第3页; 无公开对应官方解答, 编者独立补解; 反序对采用课程根为最小标号的祖先约定}
\begin{solution}
习题\ref{mit:18212-pset3-q9}的树--停车函数双射逐个保存所需统计, 因而可以证明比多项式等式更强的结论. 若 $p(i)$ 为车 $i$ 的实际车位, 定义总位移
$A(f)=\sum_i(p(i)-f(i))$. 停车成功时 $p(i)$ 遍历 $1,\ldots,n$, 所以
$A(f)=\binom{n+1}2-\sum_i f(i)\ge0$.

在上一题的分解记号下,
\begin{equation}\label{eq:18212-inversion-recursion}
 \operatorname{inv}(T)=\operatorname{inv}(T_L)+\operatorname{inv}(T_R)+r.
\end{equation}
理由是: 分支 $B$ 与外部 $C$ 的两个点不可能互为祖先. $B$ 内除了 $v$ 以外的祖先比较, 在递增重编号后不改变大小关系, 恰对应 $T_L$ 的反序对; $v$ 在内树被改为最小的根0, 其原来的反序贡献被单独移出. 原树中 $v$ 是 $B\setminus\{v\}$ 所有点的祖先, 与 $v$ 构成反序的恰是比 $v$ 小的 $r$ 个点. 外部祖先比较由 $T_R$ 保留, 根0从不贡献反序. 因而式\eqref{eq:18212-inversion-recursion}无遗漏或重复.

同时, 两群车在各自区域内按递归顺序停车. 左侧总位移为 $A(g)$; 右侧希望位和实际位都加上 $j$, 位移仍为 $A(h)$. 最后车从希望位 $j-r$ 停入 $j$, 位移为 $r$. 所以
\[
 A(f)=A(g)+A(h)+r.
\]
空树及空停车函数两统计都是0. 对大小归纳, 两个小对象上的反序数等于位移数, 上述两式遂给出每一对对应对象都有 $\operatorname{inv}(T)=A(f)$. 双射对每个对象仅取一次, 按统计取幂并求和即得题设等式.

例如前题的三棵二阶树反序数分别为 $0,0,1$, 故 $I_2(x)=2+x$; 对应停车函数的 $3-\sum f(i)$ 也分别为 $0,0,1$. 三阶例 $\{03,13,23\}$ 有反序 $(1,3),(2,3)$, 数为2, 对应 $(2,1,1)$ 的位移为 $6-4=2$. 完整三阶多项式为 $I_3(x)=6+6x+3x^2+x^3$.
\end{solution}

\begin{exercise}\label{mit:18212-pset3-q11}
\textbf{18.212 Spring 2019, Problem Set 3, Problem 11.}
证明上一题的树反序多项式满足 $I_n(-1)=A_n$, 其中 $A_n$ 为排列 $w\in S_n$ 满足 $w_1<w_2>w_3<w_4>\cdots$ 的交错排列数.
\end{exercise}
\sourceof{官方题面 \texttt{pset3.pdf} 第3页; 无公开对应官方解答, 编者独立补解; 与18.315作业3题3、4共用本册完整证明}
\begin{solution}
这道题与习题\ref{mit:18315-hw3-q3}、\ref{mit:18315-hw3-q4}的结果恰好衔接, 需核对点数和反序约定. 把本题树的顶点标号 $0,1,\ldots,n$ 同时加1, 变成 $K_{n+1}$ 的生成树, 根为最小标号1. 标号平移保留严格大小和祖先关系, 故反序数不变. 习题\ref{mit:18315-hw3-q4}通过有可逆分割的最大邻居深度优先搜索完整证明了
\[
 I_n(x)=f_{n+1}(x)=T_{K_{n+1}}(1,x).
\]
习题\ref{mit:18315-hw3-q3}又从连通图的形式指数生成函数和交错排列的递推完整证明了
$T_{K_{n+1}}(1,-1)=A_n$. 两式结合即得 $I_n(-1)=A_n$, 这里所得值本身非负, 不需绝对值. 例如 $n=2$ 时 $2-1=1=A_2$, $n=3$ 时 $6-6+3-1=2=A_3$. 若扩展到 $n=0$, 单顶点空树和空排列均计1, 结果仍成立. 本题保留18.212的独立原题映射, 共享证明的两道18.315题在同一册内, 不将未经展开的外部结论作为替代解答.
\end{solution}

\begin{exercise}\label{mit:18212-pset3-q14}
\textbf{18.212 Spring 2019, Problem Set 3, Problem 14.}
令 $C_k=\frac1{k+1}\binom{2k}k$ 为 Catalan 数. 对 $n\ge1$, 求 $n\times n$ 矩阵 $A=(C_{i+j-1})_{i,j=1}^n$ 的行列式.
\end{exercise}
\sourceof{官方题面 \texttt{pset3.pdf} 第4页; 无公开对应官方解答, 编者独立补解; 下述路径抵消引理也在本题内证明}
\begin{solution}
答案为 $\det A=1$. 用有向路径解释每个矩阵元, 再让相交路径成对抵消.

在整数格点中只保留 $x+y$ 为偶数、$y\ge0$ 的点, 允许步 $(1,1),(1,-1)$. 取起点 $S_i=(-2i+2,0)$、终点 $T_j=(2j,0)$. 从 $S_i$ 到 $T_j$ 需 $2(i+j-1)$ 步, 且不得落到横轴以下, 所以路径数为 $C_{i+j-1}$. 这个 Catalan 计数可直接由反射验证: 长度 $2r$ 的平衡路径共 $\binom{2r}r$ 条; 曾到高度 $-1$ 的路径, 把首次到达 $-1$ 前的段关于高度 $-1$ 反射, 与从高度 $-2$ 到0的路径双射, 共 $\binom{2r}{r-1}$ 条. 两者之差为 $C_r$. 我们可把整个图限制在 $-2n+2\le x\le2n$, $0\le y\le2n-1$ 的有限区域, 所有所用路径仍在其中; $x$ 严格增加, 因而是有向无环图.

展开行列式:
\[
 \det A=\sum_{\sigma\in S_n}\operatorname{sgn}(\sigma)
   \#\{(P_1,\ldots,P_n):P_i:S_i\longrightarrow T_{\sigma(i)}\}.
\]
对一个含公共顶点的路径族, 选择 $x$ 最小的相交顶点, 如有多个取 $y$ 最小者; 如有多条路径经过它, 取路径编号字典序最小的一对. 交换这两条路径在该顶点之后的尾段. 路径仍合法, 终点排列交换这两个像, 符号变反. 在所选顶点之前的全部路径均不变, 该顶点的路径编号也不变, 所以再作同一操作回到原路径族. 这就是无固定点的反号对合, 抵消全部相交项. 因而行列式只计互不相交的路径族, 带其终点排列的符号.

不相交时终点排列必须是恒等排列. 若有 $i<k$ 但 $\sigma(i)>\sigma(k)$, 则沿横轴的顺序为
$S_k,S_i,T_{\sigma(k)},T_{\sigma(i)}$. 在 $x=S_i$ 处, $P_i$ 高度为0, $P_k$ 高度须大于0; 在 $x=T_{\sigma(k)}$ 处则 $P_k$ 高度为0, $P_i$ 须大于0. 两路径在共同整数横坐标上的高度同奇偶, 高度差为偶数, 每步改变至多2. 要从正差变负差必经过差0, 即公共顶点, 矛盾. 所以 $\sigma$ 无反序, 只能为恒等排列. 这一奇偶条件也排除了两条斜边在格点以外交叉而不共有顶点的可能.

最后证明恒等终点的不相交族只有一个. 第一条 $S_1=(0,0)$ 到 $T_1=(2,0)$ 的路必为 $UD$, 峰在 $(1,1)$. 归纳设前 $i-1$ 条路均为山形
$U^{2h-1}D^{2h-1}$ $(h=1,\ldots,i-1)$. 第 $i$ 条路在第 $i-1$ 条的起点横坐标处必须有正高度, 因为同处高度0就相交. 在共同横坐标区间内, 两路高度差不能从正变为非正而不相交, 所以第 $i$ 条始终在第 $i-1$ 条之上. 在 $x=1$ 处, 第 $i-1$ 条高度为 $2i-3$, 故第 $i$ 条同奇偶的高度至少为 $2i-1$. 但它从 $S_i$ 到 $x=1$ 只有 $2i-1$ 步, 高度最多 $2i-1$, 所以此前必须全是上步; 剩余 $2i-1$ 步为回到 $T_i$ 必须全是下步. 这唯一给出山形 $U^{2i-1}D^{2i-1}$.

这些嵌套山形确实互不相交, 终点排列的符号为正, 因而唯一未抵消的项贡献1. $n=2$ 时矩阵为 $\left(\begin{smallmatrix}1&2\\2&5\end{smallmatrix}\right)$, 行列式 $5-4=1$, 可作为下指标 $i+j-1$ 的核对.
\end{solution}
